\subsubsection{Scraping}

Para el apartado de implementación de \textit{web scraping} se ha creado un fichero de código empleado el lenguaje de programación \textbf{python}. Para esta tarea se ha utilizado la librería de \textit{web scraping} denominada \textit{Playwright} con la cual se han podido obtener diferentes datos de interés para nuestro proyecto pertenecientes a la página web UFO Stalker, la cual cuenta con una amplia base de datos la cual contiene información sobre posibles avistamientos en diferentes partes del mundo. Estos datos son subidos por los usuarios o personas que han presenciado dicho avistamiento, por lo que en algunos casos puede tratarse de falsos positivos.

Otras librerías necesarias son \textit{Requests} para poder mandar los datos obtenidos mediante \textit{web scraping} a la base de datos. La librería \textit{Json}, ya que se ha creado un fichero de esta extensión con el cual se comprueba que no se obtengan datos de páginas ya visitadas anteriormente con el fin de evitar duplicados. Otra librería es \textit{Time}, la cual se utiliza para añadir \textit{sleeps} dentro del código para evitar sobrecargar la página de la que se obtiene información. También se utiliza \textit{Datetime} para modificar el formato los tiempos que se obtienen de la página web y transformarlos a formato \textit{Unix Timestamp -s}. Finalmente, se utilizan las librerías \textit{Re} y \textit{Logging}, la primera de ellas se utiliza para modificar algunos de los datos que se obtienen eliminando cualquier elementos del string obtenido que no sea un dígito, esto se hace ya que en algunos casos, en campos que deben ser enteros aparecen cosas como \textit{"Roughly 1 metter"}, de esta forma obtenemos el 1 eliminado así el resto de elementos de la string. La otra librería se utiliza para usar \textit{loggings} en vez de \textit{prints} en el código.

Para la correcta obtención de los datos se ha tenido que estudiar el HTML de dos páginas web, la primera de ellas es una tabla en la que cada fila es un avistamiento diferente, y al hacer click en cualquiera de las filas se redirige al usuario a la propia página del avistamiento que es donde se encuentran todos los datos, este análisis de los HTML se ha realizado utilizando las herramientas de desarrolladores que ofrece el navegador web Google Chrome para obtener así las URLs necesarias para posteriormente poder ir a la página de ese avistamiento y obtener los nombres de los contenedores en los cuales se encuentran los diferentes datos del avistamiento producido que queremos obtener. Una vez analizados los HTML y obtenidas las URLs y los nombres de los contenedores, se puede pasar a escribir el código necesario para la realización del \textit{web scraping}. 

A continuación se adjuntan unas imágenes de las páginas web para mostrar cómo es la tabla previamente mencionada y como se ven los datos una vez haces click en una de las filas de la tabla. Cabe destacar que la tabla tiene como máximo 10 filas, por lo que una vez se han leído esas 10 filas se debe pasar a una siguiente página:


\svgfigure[1]{UfoStalkerHome}{Captura de la tabla Ufo Stalker.}
\svgfigure[1]{UfoStalkerSighting}{Captura de un avistamiento en Ufo Stalker.}

Como ya se ha mencionado para poder realizar el \textit{web scraping} se ha creado un archivo de código empleando el lenguaje de programación \textbf{python}. Este archivo cuenta con un total de dos funciones además de un main, y obtiene actualmente un total de 15 datos de avistamientos cada dos minutos y medio.

La primera de estas funciones recibe el nombre \textbf{\textit{get\_urls()}}, y es la encargada de recorrer las filas de la tabla previamente mencionada y obtener las URLs de cada avistamiento. Para esto, esta función recibe del main dos parámetros, el primero de estos corresponde con el número de urls que deben obtenerse por cada iteración, en este caso quince, el siguiente parámetro indica a la función de que URL debe obtener los datos. Posteriormente, se abrirá el archivo json en el cual se guardan todas las url ya visitadas con el fin de evitar generar duplicados, y se creará un nuevo browser de chromium y una página de ese browser para poder navegar con \textit{Playwright}. Una vez creado el browser y cargada la página indicada por la url, entraremos en un bucle de recopilación de urls hasta llegar el número indicado como parámetro, primero mediante un \textit{try} se tratara de localizar la ventana emergente que pide consentimiento de cookies, si se localiza se hará click en el botón de consentimiento para aceptarlas y poder navegar por la página. Esta ventana emergente solo aparecerá la primera vez que se carga la página, por lo que en las demás ocasiones no generará problemas. Una vez dado el consentimiento, cerrándose así la ventana emergente, entraremos dentro del \textit{while} a un bucle \textit{for}, el cual recorrerá todas las filas de la tabla obteniendo las URLs de cada avistamiento, por cada URL de cada fila que consiga se comprobará si existe ya una entrada con esa URL en el json, de ser así esa URL no se añadirá y el contador no avanzará, si por el contrario no existe, se añadirá tanto al json como a la lista de URLs encontradas. Una vez se han obtenido todas las URLs de las filas visibles se saldrá del bucle \textit{for}, y dentro del \textit{while}, se le dará al botón de siguiente página, y se comprobará si se han obtenido todos las URLs necesarias antes de entrar de nuevo al bucle \textit{for} de la siguiente tabla. Finalmente, se retorna una lista con todas las URLs nuevas obtenidas.

La segunda de función recibe el nombre \textbf{\textit{get\_data()}}, la cual recibe un único parámetro, en este caso la lista de URLs obtenidas por la función anterior. Esta función como su nombre indica es la encargada de obtener la información relevante de cada una de las URLs en la lista y mandarlas al \textit{backend} para que se puedan añadir a la base de datos. Al igual que en la función anterior, lo primero que se hace es crear un browser de chromium y una página de ese mismo browser para poder navegar a cada URL. Una vez creados, mediante un bucle \textit{for} se recorrerán las URLs que están dentro de la lista, obteniendo por cada URL los datos relevantes para nuestro proyecto. Una vez obtenidos los datos, se pondrán en el formato requerido por el \textit{backend} y se enviaran mediante un \textit{Post}.

Finalmente, el main de este archivo es el encargado de llamar a las dos funciones previamente comentadas, este main cuenta con un bucle \textit{while True} y un \textit{sleep} de dos minutos y medio para estar constantemente preguntando por quince nuevos datos cada dos minutos y medio, de esta forma no sobrecargamos la página de la que obtenemos información. Se han utilizado otros \textit{sleep} en cada una de las funciones para no realizar las operaciones de click y volver atras de una forma muy rápida, simulando un poco lo que podría llegar a ser una navegación normal.
